\documentclass[11pt]{article}

\usepackage[margin=1in]{geometry}
\usepackage{amsmath,amssymb}
\usepackage{booktabs}
\usepackage{tabularx}
\usepackage{enumitem}
\usepackage[hidelinks]{hyperref}

\newcommand{\rDist}{\operatorname{rDist}}
\newcommand{\cDist}{\operatorname{cDist}}
\newcommand{\MinPts}{\operatorname{MinPts}}
\newcommand{\NA}{\mathrm{NA}}

\title{Planning Notes: Literature-Grounded Labeling Priority Scores}
\author{}
\date{}

\begin{document}
\maketitle

% This is a technical planning document, not a formal research paper. It
% defines the score contract, notation, formulas, provenance, and validation
% requirements for deterministic recommendation categories.

\section{Proposed Categories}

The proposed V1 retains five core categories. Each category answers a distinct
question about why a human label may be useful. Each category therefore has its
own atomic score; the five scores are not combined into a weighted sum.

\begin{table}[h]
\centering
\small
\begin{tabularx}{\textwidth}{p{0.22\textwidth} p{0.22\textwidth} X}
\toprule
Category & Labeling value & Main question \\
\midrule
Cluster Assignment Conflict & Ambiguity &
Do nearby records support the current cluster assignment? \\

Label Coverage Gap & Coverage &
Is this record in a well-supported region that has little human supervision? \\

Weak Group Reachability & Novelty / exploration &
Is this record difficult to connect to every currently known normal group? \\

Local Sparsity & Rarity / isolation &
Is this record in an unusually sparse local region? \\

Known-Outlier Similarity & Human-guided anomaly search &
Does this record resemble an anomaly that a user has already confirmed? \\
\bottomrule
\end{tabularx}
\end{table}

An optional sixth category, \emph{Model Instability}, is retained for later
evaluation. It should not be enabled in V1 until the perturbation and cluster
alignment protocol is fixed.

\section{Shared Notation and Score Contract}

\subsection{Round-Level Sets}

All quantities are computed at active-learning round $t$. The superscript
$(t)$ is retained where the round matters and may be omitted in explanatory
text when no ambiguity is possible.

\begin{description}[leftmargin=3.2cm,style=nextline]
    \item[$D$]
    The complete set of records in the current immutable dataset version.

    \item[$U_t\subseteq D$]
    Records eligible to be recommended at round $t$. This normally contains
    unlabeled records. A previously labeled record may be included only when a
    deterministic recheck reason, such as a changed assignment or a label
    conflict, has been recorded.

    \item[$D_L^{(t)}\subseteq D$]
    Records with at least one active, decisive human-provided label. A record
    marked only as \emph{uncertain} is excluded because it does not yet provide
    a semantic, normal, or outlier reference for coverage calculations.

    \item[$D_N^{(t)}\subseteq D_L^{(t)}$]
    Records that a human has confirmed as normal reference records.

    \item[$D_O^{(t)}\subseteq D_L^{(t)}$]
    Records that a human has confirmed as outliers.

    \item[$D_C^{(t)}\subseteq D$]
    Records currently assigned by SSDBCODI to a normal cluster. Records whose
    current model state is outlier or unassigned are not members of this set.

    \item[$C_{\mathrm{core}}^{(t)}\subseteq D_C^{(t)}$]
    The deterministic reliable-normal set emitted by the current SSDBCODI
    provider. In V1, this consists of records that the provider explicitly
    designates as core, reliably clustered normal records. If the provider does
    not expose this set, the Label Coverage Gap score is unavailable rather
    than reconstructed using an undocumented heuristic.

    \item[$c_t(q)$]
    The current normal-cluster identifier assigned to record $q$. It is defined
    only when $q\in D_C^{(t)}$. Cluster identifiers must be aligned through the
    session's cluster-lineage mechanism before they are compared across rounds.
\end{description}

Human semantic labels and model cluster identifiers are different objects.
Human labels describe user meaning and remain stable across rounds; $c_t(q)$
describes the current model result and may change.

Throughout the document, $|X|$ denotes the number of elements in set $X$,
$\varnothing$ denotes the empty set, and $\mathbf{1}[B]$ is $1$ when condition
$B$ is true and $0$ otherwise.

\subsection{Distance and Reachability}

Let $d_M(x,y)\geq 0$ denote the distance between records $x$ and $y$ in the
fixed model-feature space of the current dataset and preprocessing version.
For numeric data this may be Euclidean distance after the declared scaling
pipeline. For mixed tabular data, the preprocessing and distance definition
must be stored with the dataset version. Raw and transformed feature spaces
must not be mixed within one round.

Let $\MinPts\geq 1$ be the neighborhood-size parameter used by SSDBCODI. The
core distance of $q$ is
\begin{equation}
    \cDist_t(q)
    =
    d_M\!\left(q,\operatorname{NN}^{M}_{\MinPts,t}(q)\right),
\end{equation}
where $\operatorname{NN}^{M}_{i,t}(q)$ is the $i$th nearest record to $q$ under
$d_M$, excluding $q$ itself. Equal distances are ordered by stable point
identifier. The pairwise reachability distance is
\begin{equation}
    \rDist_t(x,y)
    =
    \max\!\left\{
        \cDist_t(x),
        \cDist_t(y),
        d_M(x,y)
    \right\}.
\end{equation}
These definitions require $|D|-1\geq\MinPts$. If this condition fails, every
category that depends on reachability is unavailable.

\subsection{Availability and Interpretation}

For category $j$ and candidate $q\in U_t$, the system returns
\begin{equation}
    \left(A_j^{(t)}(q),S_j^{(t)}(q),R_j^{(t)}(q)\right),
\end{equation}
where
\begin{itemize}[leftmargin=*]
    \item $A_j^{(t)}(q)\in\{0,1\}$ indicates whether the score can be computed;
    \item $S_j^{(t)}(q)\in[0,1]\cup\{\NA\}$ is the numeric score; and
    \item $R_j^{(t)}(q)$ is a deterministic explanation of availability or
    unavailability.
\end{itemize}

When a required reference set is empty, the output is $(0,\NA,R)$ rather than
an invented numeric value. Thus every eligible record receives a complete
category result even though every category is not necessarily numerically
available in every round.

A score is a deterministic ranking measure, not a calibrated probability. For
example, $S(q)=0.8$ does not mean that $q$ has an 80\% probability of being
incorrect or anomalous. Scores are compared within the same category, dataset
version, preprocessing version, and round. They must not be directly compared
across categories as though they shared a common statistical scale.

Ties are broken deterministically by stable point identifier. Batch diversity,
history penalties, and ``do not repeat in consecutive rounds'' constraints are
applied after the atomic category score has been computed. They do not change
the definition of the score itself. Ground-truth columns used for offline
evaluation must never enter any score calculation.

\subsection{Evidence Terminology}

Each category distinguishes three levels of support:
\begin{description}[leftmargin=3.8cm,style=nextline]
    \item[Formula provenance]
    Whether the formula is copied from a cited method, transformed from a cited
    formula, adapted to the current clustering setting, or newly defined for
    this dashboard.

    \item[Signal evidence]
    Whether prior work supports the measured property, such as local conflict,
    coverage, reachability, sparsity, or similarity to confirmed anomalies.

    \item[Query-policy evidence]
    Whether prior work directly shows that labeling points with high values of
    this exact score improves the intended active-learning outcome.
\end{description}

Strong signal evidence does not by itself validate a stand-alone labeling
policy. Every proposed score therefore remains subject to category-specific
query-policy evaluation.

\section{Cluster Assignment Conflict}

\paragraph{Question.}
Do nearby, structurally similar records support the current cluster assignment
of record $q$?

\paragraph{Availability.}
The score is available only if
\begin{equation}
    q\in U_t\cap D_C^{(t)}
    \quad\text{and}\quad
    |D_C^{(t)}\setminus\{q\}|\geq 1.
\end{equation}
The first condition ensures that $c_t(q)$ exists. The second ensures that at
least one comparison record exists.

\paragraph{Formula.}
Define
\begin{equation}
    k_q^{(t)}
    =
    \min\!\left(
        \MinPts,
        |D_C^{(t)}\setminus\{q\}|
    \right).
\end{equation}
Let $\mathcal{N}^{r}_{k_q,t}(q)$ be the $k_q^{(t)}$ records in
$D_C^{(t)}\setminus\{q\}$ with the smallest reachability distance to $q$. The
score is
\begin{equation}
    \boxed{
    S_{\mathrm{conflict}}^{(t)}(q)
    =
    \frac{1}{k_q^{(t)}}
    \sum_{p\in\mathcal{N}^{r}_{k_q,t}(q)}
    \mathbf{1}\!\left[c_t(p)\neq c_t(q)\right]
    }.
\end{equation}
\paragraph{Interpretation.}
$S_{\mathrm{conflict}}^{(t)}(q)$ is the fraction of nearby clustered records
that currently belong to another cluster. A value near $0$ means the local
neighborhood largely agrees with the current assignment. A value near $1$
means the neighborhood largely disagrees. A high value identifies an ambiguous
local assignment; it does not prove that the current assignment is wrong.

\paragraph{Why a label may help.}
A human label can determine whether the point belongs with its current group,
resembles a neighboring group, or is a meaningful exception. It is therefore a
useful boundary or assignment check.

\paragraph{Formula provenance.}
The score is a project-specific adaptation of $k$-Disagreeing Neighbors (kDN),
which measures the proportion of nearest neighbors with a different known class
label \cite{smith2014instance}. The adaptation replaces ground-truth class
labels with current cluster assignments and replaces ordinary nearest neighbors
with reachability neighbors. It must not be described as the original kDN
formula.

\paragraph{Literature support.}
ACCESS provides direct conceptual support for querying records near close or
overlapping cluster boundaries because constraints on these records can resolve
cluster ambiguity \cite{xu2005access}. Xiong et al. provide additional support
for active clustering based on expected uncertainty reduction
\cite{xiong2017active}. Neither paper validates this exact reachability-neighbor
fraction as a query policy.

\paragraph{Primary limitation.}
The unweighted fraction treats the nearest conflicting record and the
$k_q^{(t)}$th conflicting record equally. A reachability-weighted version may
be tested as an ablation, but it should not replace this simple baseline without
empirical evidence.

\section{Label Coverage Gap}

\paragraph{Question.}
Is record $q$ in a well-supported part of the data that has little human
supervision?

\paragraph{Availability.}
This category requires at least one active human-labeled record and restricts
the ranking to currently well-supported normal structure. Define the fixed
round-level coverage set
\begin{equation}
    E_{\mathrm{coverage}}^{(t)}
    =
    U_t\cap C_{\mathrm{core}}^{(t)}.
\end{equation}
The score is available for $q$ only when
\begin{equation}
    D_L^{(t)}\neq\varnothing,
    \qquad
    q\in E_{\mathrm{coverage}}^{(t)},
    \qquad
    |E_{\mathrm{coverage}}^{(t)}|\geq 1.
\end{equation}
The core-record gate prevents isolated noise from being interpreted as an
important uncovered region. The set $E_{\mathrm{coverage}}^{(t)}$ is frozen for
the duration of round $t$ so that filtering the displayed candidate list cannot
change the score.

\paragraph{Formula.}
For $u\in E_{\mathrm{coverage}}^{(t)}$, define its distance to the closest
human-labeled record as
\begin{equation}
    d_L^{(t)}(u)
    =
    \min_{\ell\in D_L^{(t)}}\rDist_t(u,\ell).
\end{equation}
The empirical cumulative distribution of these distances is
\begin{equation}
    \widehat F_{L,t}(z)
    =
    \frac{1}{|E_{\mathrm{coverage}}^{(t)}|}
    \sum_{u\in E_{\mathrm{coverage}}^{(t)}}
    \mathbf{1}\!\left[d_L^{(t)}(u)\leq z\right].
\end{equation}
The coverage score is the within-round percentile
\begin{equation}
    \boxed{
    S_{\mathrm{coverage}}^{(t)}(q)
    =
    \widehat F_{L,t}\!\left(d_L^{(t)}(q)\right)
    }.
\end{equation}
Tied distances receive the same score; stable point ID is used only when the
system must order tied candidates.

\paragraph{Interpretation.}
A high score means that $q$ is farther from existing human supervision than
most other well-supported candidates in the same round. It does not mean that
$q$ is anomalous. Because $q$ must pass the core-record gate, the intended case
is a plausible data region that is represented by the model but not yet
represented by human labels.

\paragraph{Why a label may help.}
A label on a representative record from an under-covered region can establish
a new semantic reference and may reduce the need to label many nearby records.
This category prevents the system from spending its entire budget only on
ambiguous or unusual cases.

\paragraph{Formula provenance.}
The nearest-label percentile is a project-defined coverage score. It is not a
formula from QUIRE or core-set active learning. Its nearest-reference geometry
is motivated by farthest-first exploration and core-set selection, while the
percentile transformation makes its within-round interpretation independent of
the raw distance unit.

\paragraph{Literature support.}
Basu et al. use farthest-first exploration to obtain supervision for previously
unrepresented cluster neighborhoods \cite{basu2004active}. Core-set active
learning formalizes label acquisition as selecting a geometrically covering
subset of the data \cite{sener2018coreset}. QUIRE shows that representative and
informative query criteria can be complementary \cite{huang2014quire}. These
works support coverage as a query objective, but none validates the exact
percentile score and SSDBCODI core gate used here.

\paragraph{Primary limitation.}
The result depends on the model-feature distance and the quality of the current
core assignments. Coverage should therefore be evaluated against graph-based
coverage and cluster-stratified baselines.

\section{Weak Group Reachability}

\paragraph{Question.}
Is record $q$ difficult to connect to every currently known normal group?

\paragraph{Availability.}
The category requires at least one human-confirmed normal reference:
\begin{equation}
    D_N^{(t)}\neq\varnothing.
\end{equation}
It also requires the SSDBCODI reachability expansion structure for the current
round. If the implementation does not compute the pathwise quantity defined
below, it must report a different score name rather than silently substituting
a direct pairwise distance.

\paragraph{Pathwise reachability.}
For each confirmed normal reference $s\in D_N^{(t)}$, let $T_{t,s}$ denote the
deterministic reachability expansion structure produced by SSDBCODI when the
expansion is rooted at $s$. For a record $q$ reached from $s$, let
$\operatorname{path}_{T_{t,s}}(s,q)$ be the recorded expansion path and define
\begin{equation}
    E_{\max,t}(q,s)
    =
    \max_{e\in\operatorname{path}_{T_{t,s}}(s,q)} w_t(e),
\end{equation}
where $w_t(e)=\rDist_t(x,y)$ for an expansion edge $e=(x,y)$. If the expansion
rooted at $s$ does not reach $q$, set $E_{\max,t}(q,s)=+\infty$.

The minimum pathwise barrier from $q$ to any confirmed normal reference is
\begin{equation}
    R_t(q)
    =
    \min_{s\in D_N^{(t)}}E_{\max,t}(q,s).
\end{equation}
SSDBCODI's reachability score is
\begin{equation}
    rScore_t(q)=\exp[-R_t(q)].
\end{equation}
The labeling-priority direction is
\begin{equation}
    \boxed{
    S_{\mathrm{far}}^{(t)}(q)
    =
    1-rScore_t(q)
    =
    1-\exp[-R_t(q)]
    }.
\end{equation}
The convention $\exp(-\infty)=0$ gives disconnected records a score of $1$.

\paragraph{Interpretation.}
A high score means that every route from $q$ to a confirmed normal reference
contains a relatively weak reachability connection. This measures separation
from known normal structure. It does not distinguish by itself among a new
group, an underrepresented group, a true outlier, or an unsupported assignment.

\paragraph{Why a label may help.}
A human label can reveal whether the record extends an existing normal type,
belongs to a previously unrepresented type, or should be treated as unusual.
The category therefore supports exploration away from known normal groups.

\paragraph{Formula provenance.}
$rScore_t(q)$ and the reversed term $1-rScore_t(q)$ follow the reachability
component used by SSDBCODI \cite{deng2022ssdbcodi}. Using this component alone
to rank labeling candidates is a project-defined acquisition policy. It is not
the complete SSDBCODI outlier score, which combines several signals.

\paragraph{Literature support.}
SSDBCODI provides strong evidence that pathwise reachability to labeled normal
objects is useful for semi-supervised outlier analysis. Active constrained
clustering provides broader support for spending supervision on structure that
is not yet represented by known neighborhoods \cite{basu2004active}. Evidence
for this exact stand-alone query policy remains to be established.

\paragraph{Primary limitation.}
The exponential mapping is monotone but scale-dependent; it is not a calibrated
normalization across datasets. The score is therefore used only for within-round
ranking under a fixed preprocessing version. Cross-dataset comparisons should
use rank statistics rather than raw score values.

\section{Local Sparsity}

\paragraph{Question.}
Does record $q$ occupy an unusually sparse local part of the data?

\paragraph{Availability.}
The score is available when
\begin{equation}
    q\in U_t
    \quad\text{and}\quad
    |D\setminus\{q\}|\geq\MinPts.
\end{equation}

\paragraph{Formula.}
Let $\operatorname{NN}^{r}_{i,t}(q)$ be the $i$th nearest record to $q$ under
$\rDist_t$, excluding $q$. Define the mean local reachability distance
\begin{equation}
    LD_t(q)
    =
    \frac{1}{\MinPts}
    \sum_{i=1}^{\MinPts}
    \rDist_t\!\left(\operatorname{NN}^{r}_{i,t}(q),q\right).
\end{equation}
SSDBCODI transforms this quantity as
\begin{equation}
    lScore_t(q)=\exp[-LD_t(q)].
\end{equation}
The local-sparsity ranking score is
\begin{equation}
    \boxed{
    S_{\mathrm{sparse}}^{(t)}(q)
    =
    1-lScore_t(q)
    =
    1-\exp[-LD_t(q)]
    }.
\end{equation}

\paragraph{Interpretation.}
A high score means that the nearest reachability neighbors of $q$ are, on
average, relatively far away. The score measures \emph{absolute local sparsity}.
It does not compare the density around $q$ with the normal density of $q$'s own
cluster and is therefore not a relative-density score such as LOF.

This category differs from Weak Group Reachability. A record may be far from all
human-confirmed normal references but still lie in a dense, coherent pocket.
Conversely, a record may be locally isolated while remaining connected to an
already known group.

\paragraph{Why a label may help.}
A human label can distinguish a legitimate rare case, a small group, a true
outlier, or a possible data-quality problem. The score only identifies the
sparse setting; it does not decide which interpretation is correct.

\paragraph{Formula provenance.}
$LD_t(q)$, $lScore_t(q)$, and the reversed term $1-lScore_t(q)$ follow the local
density component used by SSDBCODI \cite{deng2022ssdbcodi}. Using that component
alone as a labeling-priority score is a project-defined acquisition policy.

\paragraph{Literature support.}
SSDBCODI provides direct support for mean reachability distance as one of its
outlier-analysis signals. LOF provides broader evidence that local density
relative to neighboring density is useful for detecting local outliers
\cite{breunig2000lof}; however, LOF does not validate this absolute-sparsity
formula. The two approaches should be compared empirically rather than treated
as equivalent.

\paragraph{Primary limitation.}
Datasets may contain valid groups with naturally different densities. An
absolute sparsity score may repeatedly prioritize a legitimate low-density
group. A relative-density alternative should therefore be included in the
ablation study.

\section{Known-Outlier Similarity}

\paragraph{Question.}
Does record $q$ resemble an anomaly that a user has already confirmed?

\paragraph{Availability.}
The score is available only when
\begin{equation}
    q\in U_t
    \quad\text{and}\quad
    D_O^{(t)}\neq\varnothing.
\end{equation}
Before the first human-confirmed outlier exists, the category returns
$(0,\NA,\text{``no confirmed outlier reference''})$ for every eligible record.

\paragraph{Formula.}
Define the distance from $q$ to its nearest confirmed outlier as
\begin{equation}
    O_t(q)
    =
    \min_{o\in D_O^{(t)}}d_M(q,o).
\end{equation}
The score is
\begin{equation}
    \boxed{
    S_{\mathrm{known\text{-}outlier}}^{(t)}(q)
    =
    \exp[-O_t(q)]
    }.
\end{equation}

\paragraph{Interpretation.}
A high score means that $q$ is close in model-feature space to at least one
human-confirmed outlier. It does not mean that $q$ is certainly anomalous. Two
nearby records can have different real-world meanings, and a single confirmed
outlier may not represent a repeatable anomaly family.

\paragraph{Why a label may help.}
This category follows up on anomaly patterns that a user has already indicated
are relevant. Labeling a high-scoring record can confirm that the pattern
repeats or show that the original outlier was an isolated case.

\paragraph{Formula provenance.}
SSDBCODI defines $simScore(q)=\exp[-dist(q,o_{\min})]$ using Euclidean distance
to the nearest labeled outlier \cite{deng2022ssdbcodi}. The formula above is
exactly the same form when $d_M$ is that Euclidean distance. If the dashboard
uses a mixed-tabular distance, it is a declared generalization. Using the score
alone as a query policy is project-defined.

\paragraph{Literature support.}
SSDBCODI directly supports similarity to a labeled outlier as a
semi-supervised anomaly signal. Active Anomaly Discovery supports the broader
principle that expert feedback should guide subsequent anomaly search because
statistical outlierness and user-relevant anomalousness are not identical
\cite{das2016aad}. AAD does not validate nearest-confirmed-outlier similarity as
this exact stand-alone policy.

\paragraph{Primary limitation.}
The top-ranked records may be near-duplicates around one confirmed outlier.
Reference provenance should therefore record which confirmed outlier produced
the minimum distance. Batch diversity is handled after scoring so the atomic
score remains unchanged.

\section{Optional Sixth Category: Model Instability}

\paragraph{Question.}
Would plausible reruns of the clustering analysis disagree about the assignment
of record $q$?

\paragraph{Status.}
This category is not part of V1. It may be enabled only after the model
perturbation protocol, out-of-sample assignment procedure, and cluster-identity
alignment have been fixed and stored in the analysis configuration.

\paragraph{Required protocol.}
Run $M\geq 2$ predefined model variants. Each variant must assign every eligible
record, including records omitted from a bootstrap sample. For the baseline
formula below, all runs must use a common aligned state set
\begin{equation}
    \mathcal{C}_t
    =
    \{1,\ldots,K_t\}\cup\{\text{outlier}\},
\end{equation}
where normal-cluster identities have been aligned to the reference run by a
declared maximum-overlap procedure. The formula is unavailable when
$|\mathcal{C}_t|\leq 1$ or when reliable alignment is impossible.

Let $c_{m,t}(q)\in\mathcal{C}_t$ be the aligned state assigned to $q$ by model
variant $m$. Define
\begin{equation}
    P_t(c\mid q)
    =
    \frac{1}{M}
    \sum_{m=1}^{M}
    \mathbf{1}\!\left[c_{m,t}(q)=c\right].
\end{equation}
Using the convention $0\log 0=0$, the normalized assignment entropy is
\begin{equation}
    \boxed{
    S_{\mathrm{instability}}^{(t)}(q)
    =
    -\frac{
        \sum_{c\in\mathcal{C}_t}
        P_t(c\mid q)\log P_t(c\mid q)
    }{
        \log|\mathcal{C}_t|
    }
    }.
\end{equation}

\paragraph{Interpretation.}
A value near $0$ means that the predefined model variants assign $q$
consistently. A high value means that plausible variants disagree. The score
measures sensitivity to the chosen perturbation protocol, not all possible
forms of uncertainty.

\paragraph{Formula provenance.}
The normalized entropy is a project-specific clustering adaptation of
committee disagreement. Query-by-Committee supports selecting examples on
which plausible models disagree \cite{seung1992qbc}. Xiong et al. provide
clustering-specific support for selecting queries according to expected
uncertainty reduction and cluster-assignment entropy \cite{xiong2017active}.
Neither source specifies this bootstrap-and-alignment formula.

\paragraph{Primary limitation.}
Cluster labels are permutation-invariant, and splits or merges may make
one-to-one alignment invalid. If variable numbers of clusters are required, a
permutation-invariant co-assignment formulation should be designed and tested
instead of forcing unreliable label alignment.

\section{Relationship to Decision-Tree Rules}

Decision-tree rules remain an explanation layer in V1. They describe the
current SSDBCODI cluster and outlier assignments but do not create those
assignments and do not modify the five atomic recommendation scores.

For a point selected by one of the five categories, the Rule Panel may display
deterministic supporting facts such as
\begin{itemize}[leftmargin=*]
    \item whether the point satisfies its current group's surrogate rule;
    \item which raw feature condition is closest to changing; and
    \item whether nearby comparison records satisfy the same rule.
\end{itemize}
These facts help explain the selected point, but rule agreement must not be
presented as proof of semantic correctness. If a rule-based acquisition score
is added later, it should be introduced and validated as a separate atomic
category rather than hidden inside the existing scores.

\section{Validation Requirements}

\subsection{Determinism and Completeness}

For a fixed dataset version, preprocessing version, round state, and parameter
configuration:
\begin{enumerate}[leftmargin=*]
    \item every $q\in U_t$ must receive an availability flag, score or $\NA$,
    and deterministic reason for every category;
    \item repeated execution must produce identical scores, ties, and point
    order; and
    \item changing a displayed candidate filter must not alter the frozen
    round-level reference sets or score values.
\end{enumerate}

\subsection{Category Distinctness}

For every pair of available categories $a$ and $b$, measure both score
correlation and top-$B$ recommendation overlap. Define their common numeric
domain as
\begin{equation}
    Q_{a,b}^{(t)}
    =
    \{q\in U_t:A_a^{(t)}(q)=A_b^{(t)}(q)=1\}.
\end{equation}
Both statistics are unavailable when the common domain is too small for the
requested calculation. Otherwise,
\begin{equation}
    \rho_{a,b}
    =
    \operatorname{Spearman}_{q\in Q_{a,b}^{(t)}}
    \!\left(S_a^{(t)}(q),S_b^{(t)}(q)\right),
\end{equation}
\begin{equation}
    J_{a,b}^{(B)}
    =
    \frac{|\operatorname{Top}_B(S_a)\cap\operatorname{Top}_B(S_b)|}
    {|\operatorname{Top}_B(S_a)\cup\operatorname{Top}_B(S_b)|}.
\end{equation}
Here $\operatorname{Top}_B(S_j)$ contains the highest-ranked
$\min(B,|Q_{a,b}^{(t)}|)$ records from the common domain, with ties resolved by
stable point identifier.
High correlation alone does not require merging two categories if they answer
different user questions, but consistently high correlation and high top-$B$
overlap indicate that the product distinction may not justify its complexity.

\subsection{Query-Policy Effectiveness}

Offline evaluation may use hidden ground truth only as an oracle for measuring
results. It must not enter score computation. Compare each category against
random selection and relevant simple baselines using:
\begin{itemize}[leftmargin=*]
    \item label efficiency over multiple active-learning rounds;
    \item reduction in unresolved assignment conflicts;
    \item discovery of previously uncovered semantic classes;
    \item discovery rate of user-relevant anomalies; and
    \item stability of selected points across repeated identical runs.
\end{itemize}

The V1 category set should be retained only if each category contributes useful
cases that are not already supplied by the others. Model Instability should be
added only if its additional value justifies the repeated-analysis cost.

\begin{thebibliography}{99}

\bibitem{basu2004active}
S. Basu, A. Banerjee, and R. J. Mooney,
``Active Semi-Supervision for Pairwise Constrained Clustering,''
in \emph{Proceedings of the 2004 SIAM International Conference on Data Mining},
pp. 333--344, 2004.
\url{https://doi.org/10.1137/1.9781611972740.31}

\bibitem{breunig2000lof}
M. M. Breunig, H.-P. Kriegel, R. T. Ng, and J. Sander,
``LOF: Identifying Density-Based Local Outliers,''
in \emph{Proceedings of the ACM SIGMOD International Conference on Management
of Data}, pp. 93--104, 2000.
\url{https://doi.org/10.1145/342009.335388}

\bibitem{das2016aad}
S. Das, W.-K. Wong, T. G. Dietterich, A. Fern, and A. Emmott,
``Incorporating Expert Feedback into Active Anomaly Discovery,''
in \emph{Proceedings of the IEEE International Conference on Data Mining},
pp. 853--858, 2016.
\url{https://doi.org/10.1109/ICDM.2016.0102}

\bibitem{deng2022ssdbcodi}
J. Deng and E. T. Brown,
``SSDBCODI: Semi-Supervised Density-Based Clustering with Outliers Detection
Integrated,''
arXiv:2208.05561, 2022.
\url{https://arxiv.org/abs/2208.05561}

\bibitem{huang2014quire}
S.-J. Huang, R. Jin, and Z.-H. Zhou,
``Active Learning by Querying Informative and Representative Examples,''
\emph{IEEE Transactions on Pattern Analysis and Machine Intelligence},
vol. 36, no. 10, pp. 1936--1949, 2014.
\url{https://doi.org/10.1109/TPAMI.2014.2307881}

\bibitem{sener2018coreset}
O. Sener and S. Savarese,
``Active Learning for Convolutional Neural Networks: A Core-Set Approach,''
in \emph{International Conference on Learning Representations}, 2018.
\url{https://openreview.net/forum?id=H1aIuk-RW}

\bibitem{seung1992qbc}
H. S. Seung, M. Opper, and H. Sompolinsky,
``Query by Committee,''
in \emph{Proceedings of the Fifth Annual ACM Workshop on Computational Learning
Theory}, pp. 287--294, 1992.
\url{https://doi.org/10.1145/130385.130417}

\bibitem{smith2014instance}
M. R. Smith, T. R. Martinez, and C. Giraud-Carrier,
``An Instance Level Analysis of Data Complexity,''
\emph{Machine Learning}, vol. 95, no. 2, pp. 225--256, 2014.
\url{https://doi.org/10.1007/s10994-013-5422-z}

\bibitem{xiong2017active}
C. Xiong, D. M. Johnson, and J. J. Corso,
``Active Clustering with Model-Based Uncertainty Reduction,''
\emph{IEEE Transactions on Pattern Analysis and Machine Intelligence},
vol. 39, no. 1, pp. 5--17, 2017.
\url{https://doi.org/10.1109/TPAMI.2016.2539965}

\bibitem{xu2005access}
Q. Xu, M. desJardins, and K. L. Wagstaff,
``Active Constrained Clustering by Examining Spectral Eigenvectors,''
in \emph{Proceedings of the Eighth International Conference on Discovery
Science}, Lecture Notes in Computer Science, vol. 3735, pp. 294--307, 2005.
\url{https://doi.org/10.1007/11563983_25}

\end{thebibliography}

\end{document}
